% !PRINTED 262
\weberpage{262}
Ist aber $n$ eine Primzahl, so genügt der Nachweis, daß $\Phi_{1,0,n}$
keine negativen Potenzen von $q$ enthält, da das Gleiche durch
Vertauschung von $\omega$ mit $\omega:n$ für $\Phi_{n,0,1}$ folgt.

Der konstante Wert, den die Funktion $\Phi_{a,c,\partial}$ erhält, bestimmt
sich aus \gesperrt{einem} Gliede der Entwickelung.

Bei der Ausführung dieser Rechnungen dienen die Entwickelungen § 24, (11),
\weq{14}{\begin{aligned}
A &= q^{-\frac{n-1}{24}r} \prod_{1,\infty}^{h} \left(\frac{1+q^{n(2h-1)}}{1+q^{2h-1}}\right)^r \\
&\quad + q^{\frac{n-1}{24}r} \prod_{1,\infty}^{h} \left(\frac{1+q^{2h-1}}{1+q^{n(2h-1)}}\right)^r \\
B &= q^{-\frac{n+1}{24}s} \prod_{1,\infty}^{h} (1+q^{2h-1})^s (1+q^{n(2h-1)})^s \\
&\quad + \left(\frac{2}{n}\right)^{r+s} q^{\frac{n+1}{24}s} 2^s \prod_{1,\infty}^{h} \frac{1}{(1+q^{2h-1})^s (1+q^{n(2h-1)})^s} .
\end{aligned}}
Die Formeln werden nicht immer am einfachsten, wenn $r$, $s$
möglichst klein angenommen werden, sondern es erweist sich am
zweckmäßigsten, noch die Bedingung
\weq{15}{\frac{(n-1)r}{12} \equiv \frac{(n+1)s}{12} \ (\mathrm{mod}\ 2)}
hinzuzufügen.

Wir erhalten so für die sieben ersten ungeraden Primzahlen
folgende Bestimmung von $A$ und $B$:
\[
\begin{aligned}
n &= 3, & A &= \left(\frac{u}{v}\right)^6 + \left(\frac{v}{u}\right)^6, & B &= (uv)^3 - \frac{8}{(uv)^3}, \\
n &= 5, & A &= \left(\frac{u}{v}\right)^3 + \left(\frac{v}{u}\right)^3, & B &= (uv)^2 - \frac{4}{(uv)^2}, \\
n &= 7, & A &= \left(\frac{u}{v}\right)^4 + \left(\frac{v}{u}\right)^4, & B &= (uv)^3 + \frac{8}{(uv)^3}, \\
n &= 11, & A &= \left(\frac{u}{v}\right)^6 + \left(\frac{v}{u}\right)^6, & B &= uv - \frac{2}{vu}, \\
n &= 13, & A &= \frac{u}{v} + \frac{v}{u}, & B &= (uv)^6 - \frac{64}{(uv)^6}, \\
n &= 17, & A &= \left(\frac{u}{v}\right)^3 + \left(\frac{v}{u}\right)^3, & B &= (uv)^4 + \frac{16}{(uv)^4}, \\
n &= 19, & A &= \left(\frac{u}{v}\right)^2 + \left(\frac{v}{u}\right)^2, & B &= (uv)^3 - \frac{8}{(uv)^3}.
\end{aligned}
\]
